% Insert after the preserved as-is observation tail. Keep its existing section.
\hyphenrules{english}
\subsection{Mission needs, requirements and effectiveness}
\label{subsection:framework-mission-needs}
\label{subsection:framework-moes}

The expanded concerns require more than evidence storage. Engineers need verifiable requirements, controlled changes, component and cost comparisons, dynamic EPS and RAM analysis, credible models, software contracts and usable verification proposals. Configuration control connects these services. The system being developed is the EPS engineering framework. EPS hardware, flight software and ground software are engineered artifacts. Physical execution, flight operation and facility authorization remain external.

The visible source lists and complete registry preserve 144 concern occurrences, 30 roots, 114 causal edges and seven legacy alias resolutions. Table~\ref{tab:framework-mission-requirements} groups their derived needs and enterprise obligations. The representative source paths in Section~\ref{section:manage-requirements-traceability} complement the full many-to-many registry. This follows OOSEM causal, mission and black-box analysis \citep[Secs.~17.3.2--17.3.3]{friedenthal2014practical}.

The status is \textit{specified design, stakeholder review and empirical verification pending}. The preserved elicitation narrative gives no approval record for these new derivations. Literature supplies mechanisms and scoped demonstrations. It supplies no measured framework benefit.

BOM means bill of materials, DUT means device under test, and FMEA means failure modes and effects analysis. SIL and HIL mean software in the loop and hardware in the loop. TP, FP and FN count true positives, false positives and false negatives.
{\small
\setlength{\tabcolsep}{4pt}
\setlength{\parskip}{0pt}
\setlength{\LTpre}{0.4\baselineskip}
\setlength{\LTpost}{0.4\baselineskip}
\renewcommand{\arraystretch}{1.06}
\begin{longtable}{@{}>{\raggedright\arraybackslash}p{\dimexpr0.14\linewidth-0.56\tabcolsep\relax}>{\raggedright\arraybackslash}p{\dimexpr0.57\linewidth-2.28\tabcolsep\relax}>{\raggedright\arraybackslash}p{\dimexpr0.29\linewidth-1.16\tabcolsep\relax}@{}}
\caption{Derived mission needs/requirements and effectiveness observations.}\label{tab:framework-mission-requirements}\\
\toprule
\textbf{Need / MR} & \textbf{Enterprise obligation} & \textbf{MOE observation} \\
\midrule
\endfirsthead
\caption[]{Derived mission needs/requirements and effectiveness observations. (continued).}\\
\toprule
\textbf{Need / MR} & \textbf{Enterprise obligation} & \textbf{MOE observation} \\
\midrule
\endhead
\bottomrule
\endfoot
MN-01\par MR-01 & The EPS engineering enterprise shall produce reviewable framework-supported EPS and component requirements whose source, observable quantities or conditions, verification method, criterion, and unresolved decisions are explicit. & MOE-01\par Authoring/update person-hours, reviewer defects and completeness. \\
MN-02\par MR-02 & The enterprise shall identify the affected evidence for each designated change and issue a successor baseline with regenerated or explicitly stale evidence and retained discrepancy status. & MOE-02\par Correct change elapsed hours and active person-hours. TP, FP, FN and critical misses on reference discrepancies. \\
MN-03\par MR-03 & The enterprise shall obtain a revision-linked BOM and financial estimate for each designated EPS candidate, with sourced prices, quantities, assumptions, uncertainty, and missing cost data visible. & MOE-03\par BOM/cost update person-hours, arithmetic differences, missing prices and uncertainty. \\
MN-04\par MR-04 & The enterprise shall justify designated early EPS decisions using time-dependent mission and component analyses with explicit interface conditions, applicable constraints, and unsupported checks visible. & MOE-04\par Decision-relevant checks, known issues found, analysis time and effort. \\
MN-05\par MR-05 & The enterprise shall distinguish predictions supported within a documented model validity domain from extrapolation and shall distinguish a predicted control-policy benefit from a demonstrated hardware benefit. & MOE-05\par Applicable decisions, prediction residuals in their own units, uncertainty, runtime and update delay. \\
MN-06\par MR-06 & The enterprise shall relate designated RAM claims to requirements, design failure dependencies, generated analysis models, input assumptions, and verification evidence or explicit gaps. & MOE-06\par RAM preparation person-hours, failure/requirement coverage and expert-reference differences. \\
MN-07\par MR-07 & The enterprise shall provide a reviewable verification proposal and evidence contract for each designated activity, with feasibility and safety decisions reserved to the responsible engineers and testing facility. & MOE-07\par Proposal/clarification person-hours, missing fields, capture quality and synchronization uncertainty. \\
MN-08\par MR-08 & The enterprise shall produce revision-linked EPS software integration contracts, framework-specific mappings, command and mode checks, and processor, memory, timing and power evidence for the declared target configuration. & MOE-08\par Interface discrepancies, integration effort, dictionary/behavior coverage and target resources. \\
MN-09\par MR-09 & The enterprise shall permit an independent participant to reconstruct each designated analysis from its controlled baseline and evaluate agreement using result-specific comparison rules fixed before execution. & MOE-09\par Successful replays relative to attempts, independent person-hours and reconstruction gaps. \\
MN-10\par MR-10 & The enterprise shall provide a controlled review package from which a reviewer can locate the designated design, decision rationale, applicable evidence, and unresolved findings. & MOE-10\par Review effort, task completion, interpretation accuracy and clarification. \\
MN-11\par MR-11 & The enterprise shall complete the declared reference workflow in two contrasting EPS cases using an identified common core, declared case adaptations, and independently reconstructable environment instructions. & MOE-11\par Adoption/adaptation person-hours, two-case completion and direct usage cost in CAD over a declared period. \\
\end{longtable}
}

All measure values are unknown. The six original measures remain in MOE-02 (correct change and discrepancy detection), MOE-09 (reproduction) and MOE-11 (adoption, adaptation and direct cost). Additional measures address every expanded concern family. Active effort is person-time. Elapsed time, unattended execution and waiting are separate. EPS acquisition estimates differ from framework usage cost. Replay agreement differs from prediction accuracy. Achieved spacecraft RAM, reduced hardware defects and replacement of physical replicas remain external empirical outcomes.

Freeze the two contrasting cases, core, tasks, reference answers, supported profiles, participants, resources, time boundaries and price date before observations. Include clean controls, failed tasks and incomplete data. A favorable burden observation requires correct task completion. Numerical targets remain open. EPS limits, model-error rules, target budgets and facility safety/timing limits require intended-use evidence and responsible agreement.

\subsection{To-be enterprise use cases}
\label{subsection:framework-to-be-use-cases}

UC01 tailors the core and justifies an EPS option. UC02 processes a coupled EPS, load or software change. UC03 interprets a failed gate or supplied observation and prepares later verification. Engineers retain decision authority. Independent reviewers reconstruct the justification and return findings. SC01--SC03 specify these threads below. Specialist tools may be integrated dependencies or externally invoked services. Each candidate declares their placement and checked contract.
